\section{The finite packet and its exact Euler correction}\label{sec:exact-packet}

This section concerns one prescribed smooth parent solution on a closed
interval $[0,T]$. All constants are allowed to depend on this parent, on $d$,
and on the prescribed seed. The higher high coefficients use the
variational history and forward-tail construction of \cite[\S3]{OAI}.
All source constants are fixed before the packet frequency is
chosen. A bound on these constants along an infinite sequence of parents is
a separate question; none is asserted in this section.

Let $U$ be the parent, let $P_U$ be its Newton pressure, and let $\Phi_t$ be
its volume-preserving flow, with $\Phi_0=\mathrm{id}$. In particle labels put
\begin{equation}\label{eq:packet-frame}
 F=D\Phi_t,\qquad M=(DU)\circ\Phi_t=F_tF^{-1},\qquad
 H=(D^2P_U)\circ\Phi_t=-F_{tt}F^{-1}.
\end{equation}
We assume that these are the actual flow and coefficient fields, with bounded
smooth spatial derivatives and continuous coefficient paths on $[0,T]$.
For the quantitative assertions we assume a Gevrey bound on the known
coefficients $F,F^{-1},F_t,F_{tt}$:
\begin{equation}\label{eq:packet-known-jets}
 \sup_{t,y}\|D_y^j C(t,y)\|\le C_C R_C^j(j!)^2,
 \qquad C\in\{F,F^{-1},F_t,F_{tt}\},\quad j\ge0.
\end{equation}
These are conditions on the current known parent, not on the child being
constructed. The positive lower bound of $F$ follows from the actual inverse.
All derivative norms in this section are operator norms unless specified
otherwise.

Fix a unit vector $m_0\in\R^d$ and set
\[
 m(t,y)=F(t,y)^{-\top}m_0.
\]
Then $m\ne0$ and $m_t=-M^{\top}m$. The phase variable $\theta$ belongs
to $\mathbb T_P=\R/P\mathbb Z$, with $P=2\pi$ when a specific periodic seed
is used. Every vector field below has all $d$ components. The orthogonal
space $m^\perp$ has dimension $d-1$; it is never replaced by a chosen
one- or two-dimensional subspace.

For physical Sobolev persistence one may use
$s_d=\lfloor d/2\rfloor+3>d/2+1$, as in the preceding section. Packet
estimates also involve the cylinder $\R^d\times\mathbb T_P$. To avoid any
implicit three-dimensional threshold we use the larger fixed base order
\begin{equation}\label{eq:packet-base-index}
 b_d=2(d+2)+1=2d+5.
\end{equation}
The cylinder embedding, product, and commutator estimates at this order
follow by Fourier series in $\theta$ and the ordinary Sobolev estimates in
$y$. All their constants depend on $d$ and $P$, not on the number of
additional derivatives.

\subsection{The mean inverse}
Let $\mathcal H=L^2_\sigma(\R^d;\R^d)$ be the closed solenoidal subspace.
Choose a nonnegative even smooth cutoff $\chi$, equal to one near zero and
supported in the unit ball, and put $\chi_\ell(y)=\chi(\ell y)$.
There is a bounded map from the homogeneous first-order tensor space to
$L^2$ defined first on compact smooth tensors by
\[
 (T_{\chi_\ell}Q)_j
   =\sum_{i=1}^d\partial_i\{\chi_\ell(Q_{ij}-Q_{ji})\},
 \qquad A_{\chi_\ell}=T_{\chi_\ell}T_{\chi_\ell}^{*}.
\]
The tensor space here is the completion for the square sum of all
$\partial_kQ_{ij}$, with $1\le i,j,k\le d$. Hardy's inequality for
$d\ge3$ controls the term $(\partial_i\chi_\ell)Q_{ij}$ and extends $T$
continuously. Its bound is uniform for $0<\ell\le1$ after dilation.
Consequently $A_{\chi_\ell}$ is positive and symmetric, its range is
solenoidal, and every output vanishes off $\operatorname{supp}\chi_\ell$.
This is a bounded operator with spatially supported outputs; compactness
as an operator on $L^2$ is neither needed nor asserted.

We record the localization inequality used below. There are constants
$C_{1,d},C_{2,d}$ and a fixed $a_d>0$ such that, if
\begin{align*}
 \langle M(0,y)z,z\rangle&\ge-B_e|z|^2
       &&\text{for }|y|\ge a_dr/\ell,\\
 \langle M(0,y)z,z\rangle&\ge-B_c|z|^2
       &&\text{for }|y|<a_dr/\ell,
\end{align*}
where $0\le r\le1/4$, then, for $z\in\mathcal H$ and
$\mathcal L\ge C_{1,d}B_c$,
\begin{equation}\label{eq:packet-boundary-lower}
 \langle M(0)z,z\rangle_{L^2}
 +\mathcal L\langle A_{\chi_\ell}z,z\rangle_{L^2}
 \ge-(B_e+C_{2,d}B_cr^d)\|z\|_2^2.
\end{equation}
Here and below constants absorb the fixed normalization of the
cutoff plateau. To obtain \eqref{eq:packet-boundary-lower}, set
$Q=T_{\chi_\ell}^{*}z$. On the plateau, subtract one half of the uncut
antisymmetric divergence of $Q$ from $z$. The difference is weakly
harmonic, by the defining Riesz identity for $Q$. Its $L^2$ mass on the
concentric ball of relative radius $r$ is bounded by $C_dr^d$ times its
mass on the plateau ball. This follows from the interior harmonic
estimate, obtained by nested cutoff energy estimates and Sobolev
embedding, and then by dilation. The other part is controlled by
$\|Q\|_{\dot H^1}^2=\langle A_{\chi_\ell}z,z\rangle$.
Combining this local bound with the exterior lower bound proves the
inequality. The argument uses the complete tensor, so no special
three-dimensional curl formula occurs.

Assume $H(t,y)\le K I$ in quadratic-form order and
\begin{equation}\label{eq:packet-mean-guard}
 K\frac{T^2}{2}+(B_e+C_{2,d}B_cr^d)T\le\frac12.
\end{equation}
For $v\in L^2(0,T;\mathcal H)$ define
\[
 Iv(t)=-\int_t^T v(s)\,ds,\qquad X_v=F Iv,\qquad
 D_Fv=\partial_t X_v=F_tIv+Fv.
\]
The mean form is
\begin{align}
 \mathfrak a(v,w)
 &=\int_0^T\langle D_Fv,D_Fw\rangle\,dt
   -\int_0^T\langle HX_v,X_w\rangle\,dt\notag\\
 &\quad+\langle (M(0)+\mathcal L A_{\chi_\ell})X_v(0),
                         X_w(0)\rangle.\label{eq:packet-mean-form}
\end{align}
Spatial integrations are included in each inner product. The form need
not be symmetric because $M(0)$ need not be symmetric. It is bounded and
coercive. Indeed $X_v(T)=0$, hence
\[
 \|X_v\|_{L^2_tL^2_y}^2\le\frac{T^2}{2}\|D_Fv\|_{L^2_tL^2_y}^2,
 \qquad \|X_v(0)\|_2^2\le T\|D_Fv\|_{L^2_tL^2_y}^2.
\]
Moreover $X_v(0)=Iv(0)\in\mathcal H$, because $F(0)=I$.
Equations \eqref{eq:packet-boundary-lower} and
\eqref{eq:packet-mean-guard} give
$\mathfrak a(v,v)\ge\frac12\|D_Fv\|_2^2$.
The actual inverse gives a polynomial conversion. If
$I=\|F^{-1}\|_\infty$ and $C_1=\|F_t\|_\infty$, then
$v=F^{-1}D_Fv-F^{-1}F_tF^{-1}X_v$, so
$\|v\|_2\le I(1+C_1IT)\|D_Fv\|_2$.
Thus the Hilbert-space coercive inverse constructs the unique $v$ with
\begin{equation}\label{eq:packet-mean-weak}
 \mathfrak a(v,w)=-\int_0^T\langle f,X_w\rangle\,dt
 \quad\text{for every }w\in L^2(0,T;\mathcal H).
\end{equation}
Set $B=Fv$.

\begin{lemma}[Strong mean and its pressure]\label{lem:packet-strong-mean}
For a forcing $f$ with continuous $L^2$ spatial derivatives of every
order, the preceding construction gives continuous paths $B$ and $B_t$
in every spatial Sobolev space and a spatially smooth scalar $q$,
whose spatial derivatives depend continuously on time locally in space,
normalized by $q(t,0)=0$, such that
\begin{equation}\label{eq:packet-mean-pde}
 B_t+MB+F^{-\top}\nabla_yq=f,
 \qquad F^{-1}B\in\mathcal H.
\end{equation}
These identities hold pointwise, with one-sided time derivatives at both
endpoints. In addition, for an actual $z_0\in\mathcal H$,
\begin{equation}\label{eq:packet-mean-initial}
 B(0)=\mathcal L A_{\chi_\ell}z_0,
 \qquad\operatorname{supp}B(0)\subset\overline B(0,\ell^{-1}).
\end{equation}
\end{lemma}
\begin{proof}
Translate the coefficients and forcing in the spatial variable. The
coercive inverse is differentiable under each translation; differentiating
its inverse identity expresses each new derivative using previously
constructed derivatives and the known coefficient derivatives. This
produces genuine $L^2$ derivatives, rather than a formal differentiation
of a selected representative. The same argument applies to the translated
cutoff operator, using its construction on the fixed homogeneous tensor
space. Every finite derivative order is therefore available.

Let $Q_t:\mathcal H\to L^2$ be $Q_tz=F(t)z$. Its Gram operator
$Q_t^*Q_t$ has a positive lower bound from $F^{-1}$. The weak equation,
$F_{tt}=-HF$, and integration by parts give
\[
 v_t=(Q_t^*Q_t)^{-1}Q_t^*\{f-2F_tv\}.
\]
First this is an $L^2$ time identity. The $H^1$ trace construction and
continuity of the coefficients then give continuous representatives of
$v$ and $v_t$, with the same spatial derivatives. The residual
$r=f-2F_tv-Fv_t$ satisfies $Q_t^*r=0$; equivalently,
$F^{\top}r$ lies in the closed $L^2$ gradient space. Its smooth
representative is curl-free; hence
\[
 q(t,y)=\int_0^1\langle F(t,sy)^{\top}r(t,sy),y\rangle\,ds
\]
satisfies $q(t,0)=0$ and $\nabla q=F^{\top}r$. This proves
\eqref{eq:packet-mean-pde}. The pointwise equations follow from equality
of continuous representatives of the same $L^2$ field. The Gram
construction and closedness give $v(t)\in\mathcal H$ at every closed
time. In particular, $B(0)=v(0)\in\mathcal H$. With $z_0=Iv(0)$,
the boundary identity in \eqref{eq:packet-mean-weak} is
\[
 P_\sigma\{-F_t(0)z_0-B(0)+(M(0)+\mathcal L A_{\chi_\ell})z_0\}=0.
\]
Since $F_t(0)=M(0)$ and both $B(0)$ and $A_{\chi_\ell}z_0$ are
solenoidal, this yields \eqref{eq:packet-mean-initial} as an $L^2$
identity. Spatial continuity gives the asserted pointwise support.
\end{proof}

All these arguments are Hilbert-space arguments. For a transverse
endpoint problem the map $Q_t$ is replaced by $F(t)R_\perp$, where
$R_\perp:\R^{d-1}\to m_0^\perp$ is a linear isometry onto the whole
orthogonal complement. Its Gram lower bound is again supplied by
$F^{-1}$. The energy, trace, and coercive-inverse arguments are unchanged;
there are $d-1$ transverse variables throughout.

\subsection{The joined high inverse and the full tensor corrector}
Fix an activation time $0<t_0<T$. Choose an isometry
$R_\perp:\R^{d-1}\to m_0^\perp$ onto the whole orthogonal complement
and put $Q=FR_\perp$. Its range is $m(t,y)^\perp$, and
$Q_t=MQ$, $Q_{tt}=-HQ$. The same fixed space $\R^{d-1}$ is used at
every label and time; it does not discard any extra transverse direction.

For an angular-mean-zero forcing
$f\in C([0,T];H^q(\R^d\times\mathbb T_P;\R^d))$ for every finite $q$,
first solve on $[0,t_0]$ the following displacement problem. Find
$z\in H^1_0(0,t_0;L^2(\R^d\times\mathbb T_P;\R^{d-1}))$ such that
\begin{equation}\label{eq:packet-high-history}
 \int_0^{t_0}\!\bigl[\langle (Qz)_t,(Q\zeta)_t\rangle
            -\langle HQz,Q\zeta\rangle\bigr]dt
       =-\int_0^{t_0}\!\langle f,Q\zeta\rangle dt
 \quad(\zeta\in H^1_0).
\end{equation}
Both displacement traces are zero. This condition does not prescribe the
initial high velocity. Set $A=Qz_t$ on the history interval.

\begin{lemma}[The bounded joined high inverse]\label{lem:packet-joined-high}
Under the known frame bounds and $Kt_0^2/2\le1/2$, problem
\eqref{eq:packet-high-history} has a unique solution. Its velocity $A$,
its time derivative, and their spatial derivatives are continuous on
$[0,t_0]$. At $t_0$ its actual velocity trace can be continued by the
forward high equation. This gives a unique joined high operator on
$[0,T]$ with the specified history boundary conditions. The history
inverse and trace costs are polynomial in the known coefficient norms,
inverse lower bound, $t_0$ and $t_0^{-1}$. No exponential forward
estimate on $[0,t_0]$ is used.
\end{lemma}
\begin{proof}
Write $X=Qz$. The one-sided Poincar\'e estimate and the curvature upper
bound give
\[
 \int_0^{t_0}(|X_t|^2-\langle HX,X\rangle)dt
       \ge\tfrac12\|X_t\|_{L^2}^2.
\]
If $Q^*Q\ge cI$, then $z=(Q^*Q)^{-1}Q^*X$, and direct differentiation,
using $X_t=Q_tz+Qz_t$, gives the polynomial estimate
\[
 \|z_t\|_2\le
 c^{-1/2}(1+t_0 c^{-1/2}\|Q_t\|_\infty)\|X_t\|_2.
\]
Thus the form is coercive on the fixed Hilbert space of coordinate
velocities $z_t$ with zero time integral. Its bounded coercive inverse
constructs $z$. Integration by parts and $Q_{tt}=-HQ$ give
\[
 (Q^*Q)z_{tt}=Q^*(f-2Q_tz_t).
\]
The actual positive Gram inverse therefore gives $z_{tt}$ and the
continuous $H^1$ traces. Repeated differentiation under spatial
translation is justified on this same fixed Hilbert space, exactly as
for the mean inverse. It gives all the stated derivatives and preserves
both homogeneous displacement boundary conditions.

At order zero, a physical-space estimate avoids the frame
conversion altogether. Put $M_* =\sup\|M\|$. Testing the form with $z$
gives $\|X_t\|_2\le2t_0\|f\|_2$. Since $A=X_t-MX$, it follows that
\[
 \|A\|_2\le2t_0(1+M_*t_0)\|f\|_2.
\]
The velocity equation derived below implies
$\|A_t\|_2\le3M_*\|A\|_2+\|f\|_2$. The elementary Hilbert-valued
trace estimate consequently gives
\begin{equation}\label{eq:packet-history-polynomial-trace}
 \|A\|_{C_tL^2}\le
 \bigl[2t_0(1+M_*t_0)(t_0^{-1/2}+3M_*\sqrt{t_0})
                         +\sqrt{t_0}\bigr]\|f\|_{L^2_tL^2}.
\end{equation}
This also bounds the initial datum for the forward segment.
\end{proof}

The Gram equation and tangency show that the same $A$ satisfies
\begin{equation}\label{eq:packet-high-ode}
 A_t=\left(-M+2\frac{m\otimes M^{\top}m}{|m|^2}\right)A
       +\left(I-\frac{m\otimes m}{|m|^2}\right)f.
\end{equation}
Here $(a\otimes b)v=a\langle b,v\rangle$. Indeed the residual
$f-A_t-MA$ is perpendicular to the whole range of $Q$, hence is a
multiple of $m$. Differentiating $m\cdot A=0$ gives its coefficient.
For $t\ge t_0$ use the actual fundamental evolution of
\eqref{eq:packet-high-ode} and Duhamel's formula with datum $A(t_0)$.
The velocities and their derivatives match at $t_0$ by the common
right-hand side, so the joined field is strongly differentiable there
as well as at the two closed endpoints.

Let $g$ be a positive continuous function with $g(t_0)=1$, and assume
the supported forward bound
\begin{equation}\label{eq:packet-tail-propagator}
 \|\mathcal U(t,s;y)\|\le C_{\rm f}\,g(t)/g(s),
 \qquad t_0\le s\le t\le T,
 \quad y\text{ in the fixed support neighborhood}.
\end{equation}
This hypothesis concerns the known parent evolution on the activation
tail. The polynomial history estimate is supplied separately by
\eqref{eq:packet-history-polynomial-trace}.

Define the normalized periodic primitive $\partial_\theta^{-1}$ and put
\begin{equation}\label{eq:packet-high-pressure}
 \pi=\partial_\theta^{-1}
       \frac{\langle m,f\rangle-2\langle m,MA\rangle}{|m|^2}.
\end{equation}
Then $A_t+MA+m\partial_\theta\pi=f$. The inverse, trace, and forward
maps commute with the angular integral. They therefore preserve zero
angular mean. They also commute with multiplication by a fixed spatial
cutoff, since these linear high equations contain no spatial derivative.
Their uniqueness proves support preservation and covariance. Thus the
joined high field and its pressure are smooth in $(y,\theta)$, periodic,
and supported where the forcing is supported. Their spatial derivatives
are continuous in time; the high velocity also has the continuous time
derivative given by \eqref{eq:packet-high-ode}.

The primary uses a nonzero terminal displacement. To distinguish its
physical and coordinate forms, let the selected physical endpoint be
$Y_{\rm phys}\in m(t_0,0)^\perp$, and fix once and for all
\[
 Y=R_\perp^*F(t_0,0)^{-1}Y_{\rm phys}\in\R^{d-1}.
\]
Then $Q(t_0,0)Y=Y_{\rm phys}$ and
$|Y|\le\|F(t_0,0)^{-1}\|\,|Y_{\rm phys}|$. This same coordinate
vector is used at every label; the endpoint is not selected again as
the label varies.

The actual cylinder problem has terminal datum
$\alpha\chi_{\rm in}(y)f_\delta(\theta)Y$, which belongs to every
cylinder Sobolev space, and initial displacement zero. Subtract the
affine lift $\alpha(t/t_0)\chi_{\rm in}f_\delta Y$. Using
$Q_{tt}=-HQ$, its remaining zero-trace problem is exactly
\eqref{eq:packet-high-history} with forcing
$-2\alpha t_0^{-1}\chi_{\rm in}f_\delta Q_tY$. This forcing is
continuous in every cylinder Sobolev space and has the prescribed
compact spatial support.

Equivalently, at each fixed label solve the finite-dimensional homogeneous
displacement problem with endpoints $0,Y$, and write $b_Y(t,y)=Qz_t$
for its velocity, continued homogeneously on the tail. The uniform frame
bounds justify this pointwise solve and its label derivatives. Linearity
and uniqueness identify the actual cylinder velocity with the compact
periodic primary
\begin{equation}\label{eq:packet-joined-primary}
 A_1(t,y,\theta)=\alpha\chi_{\rm in}(y)f_\delta(\theta)b_Y(t,y).
\end{equation}
Its initial value generally depends on $y$. The first stage,
whose activation time is zero, instead uses a prescribed homogeneous
forward datum. All subsequent algebra is identical. For higher grades
the displacement endpoints in \eqref{eq:packet-high-history} are both
zero, but their initial velocities are retained.

For every such $A$ define all $d^2$ entries of the potential
\begin{equation}\label{eq:packet-tensor-potential}
 Q_{ij}=\partial_\theta^{-1}
       \frac{m_iA_j-m_jA_i}{|m|^2},\qquad Q_{ij}=-Q_{ji},
 \qquad
 C_j=\sum_{i=1}^d D_yQ_{ij}[F^{-1}e_i].
\end{equation}
If $\kappa=k^{-1}$ and $\psi=\Phi_t^{-1}$, the chain rule gives
\begin{equation}\label{eq:packet-pair-curl}
 A_j(\psi(x),km_0\cdot\psi(x))+
 \kappa C_j(\psi(x),km_0\cdot\psi(x))
 =\kappa\sum_{i=1}^d\partial_{x_i}
    Q_{ij}(\psi(x),km_0\cdot\psi(x)).
\end{equation}
Indeed the phase part on the right is
$\sum_i m_i\partial_\theta Q_{ij}=A_j$, by tangency.
Taking divergence and interchanging the two ordinary derivatives cancels
all terms by antisymmetry. This proof retains every pair $(i,j)$,
including mixed active and extra directions. Applying the actual
volume-preserving Piola transformation shows that
$F^{-1}(A+\kappa C)$ has zero divergence on every phase-shifted label
graph. Integration over the shift gives its membership in the cylinder's
closed divergence-free $L^2$ space. The potentials are smooth and have
the same compact spatial support as $A$, so the integrations by parts
are legitimate. The time derivative of $C$ is obtained by differentiating
\eqref{eq:packet-tensor-potential}: it includes the derivatives of $F^{-1}$,
$m$, $|m|^{-2}$, the primitive, and $A$. This is the actual time derivative
at the closed endpoints as well.

\subsection{A single strict-prefix recursion}
For vector fields $a,b$ on the time-label-angle cylinder write
\[
 S(a,b)=D_yb[F^{-1}a],\qquad
 T(a,b)=\langle m,a\rangle\partial_\theta b,\qquad
 L a=\partial_ta+Ma.
\]
Let $A_1$ be the actual homogeneous joined primary
\eqref{eq:packet-joined-primary} (or the first-stage forward primary). Put $B_1=q_1=0$, form $C_1$ by
\eqref{eq:packet-tensor-potential}, and set every grade-zero field to zero.
Suppose grades below $p\ge2$ have been constructed. Set
\[
 V_i=A_i+B_i+C_{i-1}\quad(1\le i<p),\qquad
 \widehat V_p=C_{p-1},
\]
and extend this known sequence by zero above $p$. The literal known force is
\begin{equation}\label{eq:packet-known-force}
 K_p=-LC_{p-1}-F^{-\top}\nabla_y\pi_{p-1}
       -\sum_{i+j=p}S(\widehat V_i,\widehat V_j)
       -\sum_{i+j=p+1}T(\widehat V_i,\widehat V_j).
\end{equation}
Every term uses only the strict prefix. Let an overline denote normalized
angular averaging. First solve the actual mean problem with forcing
$\overline K_p$, obtaining $B_p,q_p$. Then apply the joined high
inverse of Lemma~\ref{lem:packet-joined-high} with forcing
\begin{equation}\label{eq:packet-new-high-force}
 f_p=K_p-\overline K_p-\langle m,B_p\rangle\partial_\theta A_1,
\end{equation}
obtaining $A_p,\pi_p$, and form $C_p$ from $A_p$.
The final term of \eqref{eq:packet-new-high-force} has zero angular mean.
Terms in $K_p$ containing only means are angularly constant and disappear
from $K_p-\overline K_p$. Every remaining term contains a high field or
corrector, so it has the same compact spatial support as the seed. The
actual mean forcing has all continuous $L^2$ spatial derivatives by the
Sobolev product estimates. The high forcing is smooth, periodic, mean
zero, and spatially supported. These facts inductively verify the domains
of the two solvers.

There is no simultaneous nonlinear solve at grade $p$. In the quadratic
coefficient, replacing $\widehat V_p$ by
$\widehat V_p+A_p+B_p$ adds only
$\langle m,B_p\rangle\partial_\theta A_1$.
The other potentially new fast terms vanish because $A_p,A_1$ are
tangent and $B_p$ is angularly constant. This proves both strict-prefix
coherence and the grade-$p$ momentum equation. In particular, families
constructed for different finite cutoffs are restrictions of one family.

\subsection{Finite truncation, divergence, and the residual identity}
For $N\ge1$ define
\begin{align}
 V^{[N]}&=\sum_{i=1}^N\kappa^i(A_i+B_i+\kappa C_i),\notag\\
 p^{[N]}&=\sum_{i=1}^N\kappa^i(q_i+\kappa\pi_i).
                                                        \label{eq:packet-finite-pairs}
\end{align}
Equivalently the velocity grades are $V_i=A_i+B_i+C_{i-1}$ for
$1\le i\le N$, followed by the terminal grade $V_{N+1}=C_N$.
The pressure grades are $p_i=q_i+\pi_{i-1}$ for $1\le i\le N$,
followed by $p_{N+1}=\pi_N$. Grades outside this range are zero.
In particular the terminal grade does not include unsolved
$A_{N+1}$ or $B_{N+1}$.

Each high-corrector pair in \eqref{eq:packet-finite-pairs} is exactly
solenoidal after graph evaluation by \eqref{eq:packet-pair-curl}.
Each mean has $F^{-1}B_i\in\mathcal H$ by
Lemma~\ref{lem:packet-strong-mean}; its constant-angle embedding therefore
belongs to the same cylinder solenoidal space. Thus the complete finite
sum satisfies the closed divergence constraint. The constant-angle
embedding of an ordinary $L^2$ gradient lies in the cylinder's closed
gradient space: approximate it by gradients of compactly supported
smooth scalars and use $\nabla_\kappa\varphi=\kappa\nabla_y\varphi$
for an angle-independent scalar. This applies to the mean pressure
gradients. The compact periodic high scalar pressures themselves give
gradient generators. Their finite sum therefore belongs to the required
closed gradient space.

For any raw $V,p$ define its label momentum residual by
\begin{equation}\label{eq:packet-raw-residual}
 \mathcal R_\kappa(V,p)=LV+F^{-\top}\nabla_yp+
 \kappa^{-1}m\partial_\theta p+S(V,V)+\kappa^{-1}T(V,V).
\end{equation}
For the finite grade sequences above, let
\[
 R_n=LV_n+F^{-\top}\nabla_yp_n+m\partial_\theta p_{n+1}
       +\sum_{i+j=n}S(V_i,V_j)+\sum_{i+j=n+1}T(V_i,V_j).
\]
The actual mean and high equations give $R_n=0$ for $0\le n\le N$.
Expanding the finite sums with their actual derivatives, including the
one-sided time derivatives, consequently gives the exact identity
\begin{equation}\label{eq:packet-tail-identity}
 \mathcal R_\kappa(V^{[N]},p^{[N]})
   =\sum_{n=N+1}^{2N+2}\kappa^nR_n.
\end{equation}
Neither a time derivative of $p$ nor an infinite formal expansion is used.
In particular, $R_{N+1}$ retains $LC_N$ and
$F^{-\top}\nabla_y\pi_N$.

\subsection{Quantitative bounds at one known parent}
We give the estimates in a form that makes the derivative counting
explicit. For an integer $a\ge0$ set
$\mathcal G_a(j)=R^{j+a}((j+a)!)^2$.
Use ordered cylinder derivatives with alphabet $\{1,\ldots,d+1\}$.
More precisely, the block of order $j$ is
\[
 \operatorname{Block}_j(u)=
 \sum_{|w|=j}\sum_{l=0}^{b_d}\sum_{|v|=l}
     \|\partial^v\partial^w u\|_{C([0,T];L^2)}.
\]
At fixed time the corresponding word Sobolev norm is equivalent to the
usual $H^{b_d}$ norm, with fixed constants depending only on $d,b_d$.
A positive time weight is inserted by dividing the actual field by that
weight before taking spatial or angular derivatives. For one sufficiently
large $R$, depending only on the known coefficients, seed, and solver
bounds, the preceding recursion gives
\begin{align}
 \operatorname{Block}_j(B_p,B_{p,t},\nabla_yq_p)
   &\le H_0^{2p-2}\mathcal G_{100p-140}(j),&&p\ge2,\notag\\
 \operatorname{Block}_j(A_p,A_{p,t},C_p,C_{p,t},\nabla_y\pi_p)
   &\le \gamma(t)H_0^{2p-2}\mathcal G_{100p-80}(j),&&p\ge1.
                                                        \label{eq:packet-grade-bounds}
\end{align}
Here each entry has the stated bound separately. We use
$\gamma(t)=\alpha$ for $0\le t\le t_0$ and
$\gamma(t)=\alpha g(t)$ on the tail, with
$H_0\ge\max(1,\sup\gamma)$. The notation with $\gamma(t)$ means the
uniform-time block of the field divided by $\gamma$. The supported
forward constant in \eqref{eq:packet-tail-propagator}, the history trace
cost, and the actual terminal datum are included among the known source
constants controlling $R$. In particular the higher initial high fields
are bounded with the factor $\alpha$, not set to zero. This is the
family used below and in Section~\ref{sec:initial-summability}.

For completeness, the derivative estimates underlying
\eqref{eq:packet-grade-bounds} are as follows. The word Leibniz rule is
bounded by a binomial convolution. Its Gevrey convolution closes because
\[
 \binom jl(l!)^2((j-l)!)^2
       =\frac{(j!)^2}{\binom jl},
 \qquad \sum_{l=0}^j {\binom{j}{l}}^{-1}\le3.
\]
Fixed-order Sobolev products add a constant depending on $b_d,d$.
On the history interval, differentiate the actual fixed-space coercive
inverse and then use the Gram equation and its trace. On the tail,
differentiate the actual Duhamel equation and use
\eqref{eq:packet-tail-propagator}. The history terminal trace is the
datum in this Volterra recurrence. The joined high velocity and its time
derivative consume at most three grades. Differentiating the actual coercive
inverse, followed by its Gram reconstruction and the $H^1$ time trace,
consumes at most three grades for the mean, its time derivative, and its
pressure gradient. These are estimates on word blocks with the full
alphabet. The angular mean is a bounded $L^2$ map with norm $P^{-1/2}$,
and the constant-angle embedding has norm $P^{1/2}$; the factors cancel.
Thus converting mean input to mean output does not multiply the unknown
word radius at each grade. Only the known coefficient radius is enlarged
once. Also the mean time weight is the constant $H_0^{2p-2}$, so no
commutation of a time-dependent weight with the mean inverse is claimed.
The history estimate is applied to $f/\alpha$ and the tail estimate
to $f/(\alpha g)$; linearity preserves the common factor $\alpha$.
The tail uses the weighted Duhamel operator and never differentiates $g$. The tensor potential and its actual time derivative cost one spatial
derivative and a fixed $d$-dependent multiplier constant; the pressure
source \eqref{eq:packet-high-pressure} and one slow derivative cost at
most two more grades. Enlarging $R$ once absorbs these fixed constants.
The grade counting can be made explicit. The mean force has shift at
most $100p-150$, and the high force has shift at most $100p-90$.
For example a slow high--high term with $i+j=p$ has shift
$100p-159$. In a fast term with $i+j=p+1$, the first high factor is
annihilated by $m\cdot A_i=0$. A first corrector factor has shift
$100i-180$, so the resulting high term has shift at most $100p-159$;
a first mean factor gives shift $100p-119$, and has zero angular average.
The new term $(m\cdot B_p)\partial_\theta A_1$ has the latter shift as
well. These leave the displayed force margins. The profile powers close
because $\gamma\le H_0$; in particular a mean--high fast product has
exactly the allowed factor $\gamma H_0^{2p-2}$.
The at most quadratic number of summands is absorbed by the spare
factorial shifts: for $a\ge0$,
$\mathcal G_{a+2}(j)/\mathcal G_a(j)
 =R^2(j+a+1)^2(j+a+2)^2$.
The low grades are covered by one enlargement of the known radius.
The mean inverse and the joined high inverse, followed by the corrector
and pressure estimates, use fewer than the ten reserved shifts.
Strong induction proves \eqref{eq:packet-grade-bounds} for this same
joined family, without an estimate on a yet unconstructed output.

For the correction use the truncated weighted norm
\[
 \|u\|_{r,b_d;J}
 =\sum_{j=0}^J\frac{r^j}{(j!)^2}
     \sum_{|w|=j}\sum_{l=0}^{b_d}\sum_{|v|=l}\|\partial^v\partial^wu\|_2.
\]
The final conversion from a block bound
$C\mathcal G_a(j)$ costs at most $2C(4R)^a(a!)^2$ when $4Rr\le1/2$.
This conversion is performed after the word recursion, not at every
recursive step. In particular, at every derivative cutoff,
\[
 \|V_i\|_{r,b_d;J}\le
  6H_0^{2i}(4R)^{100i}((100i)!)^2\quad(i\ge1),
\]
with analogous bounds for the actual time derivatives and pressure
forces at a smaller fixed radius. For explicit expansion-order counting, choose a known-source constant
$C_0\ge1$ and put $S_N=C_0(N+1)^{200}$. The inequality
$(100i)!\le(100(N+1))^{100i}$ for $i\le N+1$ bounds each displayed
grade by $S_N^i$ after enlarging $C_0$. The actual time derivatives
and pressure terms obey the same bound; one slow or angular derivative
is paid by one fixed radius reduction, whose reciprocal cost is included
in $C_0$. A slow product with $i+j=n$ costs at most $C S_N^n$, and a
fast product with $i+j=n+1$ at most $C S_N^{n+1}$. There are at most
$(N+2)^2$ terms. The terminal linear terms have the same bound, including
the true $C_N$ time derivative and the slow derivative of $\pi_N$.
Thus $C_1(N+2)^2S_N^{n+1}$ bounds the entire residual grade. Taking
$B_N=C_2(N+1)^{202}$ with a sufficiently large fixed $C_2$ absorbs this
factor for every $n\ge N+1$. The convenient larger exponent $300$ gives
\begin{equation}\label{eq:packet-tail-base}
 B_N\le C_*(N+1)^{300},\qquad
 \|R_n\|_{r,b_d;J}\le B_N^{n+1}\quad(N+1\le n\le2N+2).
\end{equation}
The constant $C_*$ is fixed by the current parent and seed, and is
independent of $N,k,J$. Spatial derivatives of the high pressure, the
time derivative of $C_N$, and the terminal linear terms are all included
in this estimate.

Put
\begin{equation}\label{eq:packet-coordinate-change}
 Z=kF^{-1}V^{[N]},\quad Q^a=k^2p^{[N]},\quad
 \nabla_\kappa=\kappa\nabla_y+m_0\partial_\theta,
 \quad G=F^{-1}F^{-\top}.
\end{equation}
Then the actual coordinate residual is $\mathcal E=kF^{-1}\mathcal
R_\kappa(V^{[N]},p^{[N]})$. If $B_N\le k^{1/100}$ and the known inverse
multiplier cost is at most $k^{1/100}$, the geometric sum in
\eqref{eq:packet-tail-identity} yields, for $N\ge X-1$, $X\ge6$,
\begin{equation}\label{eq:packet-small-residual}
 \|\mathcal E\|_{r,b_d;J}
       \le \exp(-\tfrac7{10}X\log k).
\end{equation}
The field $Z$ and its first cylinder derivatives have fixed weighted
bounds $B_0,B_1$. More sharply, its actual cylinder transport vector is
\[
 \mathcal V_\kappa Z=(\kappa Z,m_0\cdot Z),
 \qquad \|\mathcal V_\kappa Z\|_{r,b_d;J}\le D_*/k.
\]
The order-one high field is tangent, so its fast transport component
vanishes exactly. All other contributions carry at least one inverse
frequency. This fact is essential: using the full $O(1)$ norm of $Z$
in place of its small transport norm would lose the needed radius bound.

\subsection{An exact correction on the cylinder}
The coordinate equation is
\begin{equation}\label{eq:packet-coordinate-equation}
 W_t+(W\cdot\nabla_\kappa)W+2F^{-1}F_tW
    +\kappa\sum_{i=1}^d W_iF^{-1}(\partial_iF)W
    +G\mathbf p=0,
 \qquad \nabla_\kappa\cdot W=0.
\end{equation}
Here $\mathbf p$ is a vector pressure variable in
$\mathscr G_\kappa=\overline{\{\nabla_\kappa\varphi:
\varphi\text{ smooth, periodic, compact in }y\}}^{L^2}$.
The notation does not assert that $\mathbf p$ is the derivative of a
global periodic scalar. In particular no closed-range assertion for the
degenerate $d$-component operator on the $(d+1)$-dimensional cylinder is
used. All $d$ quadratic coefficient matrices occur in this equation.
Substituting $Z$ and the actual approximate pressure vector
$\mathbf p^a=\nabla_\kappa Q^a$ gives exactly the residual $\mathcal E$
above. The correction pressure is solved in this closed subspace of
vector-valued cylinder $L^2$.
The positive matrix $G$ gives a coercive variational pressure inverse on
that subspace. Its derivatives are controlled by the true coefficient
jets and commutators.

Seek $W=Z+E$ with $E(0)=0$. Use the inverse metric
$G^{-1}=F^{\top}F$ in the finite weighted Sobolev energy. The top
pressure term cancels since
$\langle G^{-1}E,G\mathbf p^e\rangle
 =\langle E,\mathbf p^e\rangle=0$ for $\mathbf p^e\in\mathscr G_\kappa$.
The top transport term is integrated by parts using the closed
solenoidal constraint. Differentiated metric terms and lower products
have a known coefficient bound. To specify the energy convention, put
$\|u\|_{F^{\top}F}=(\int\langle F^{\top}Fu,u\rangle)^{1/2}$ and set
\[
 \mathcal A_J=\sum_{j=0}^J\frac{r^j}{(j!)^2}
   \sum_{|w|=j}\sum_{l=0}^{b_d}\sum_{|v|=l}
       \|\partial^v\partial^w E\|_{F^{\top}F}.
\]
Thus $\mathcal A_J$ is a sum of metric norms, not a squared norm; it is
uniformly equivalent to the preceding weighted norm by the known frame
bounds. Let $\mathcal L_J$ be the same sum with the additional factor
$j$ in every summand. Derivatives of norms at zero can be justified by
first replacing each norm by its quadratic regularization and then
letting the regularization parameter decrease to zero.

The orthogonal gradient projection commutes with every cylinder
derivative. In the differentiated pressure equation, each nontrivial
commutator with $G$ lowers the derivative order on the pressure by at
least one. This compensates for the one derivative in the transport
source. After the fixed-base Sobolev pressure estimate, its remaining
coefficient-word convolutions are absorbed by choosing the known radius
small enough. Thus no unknown derivative of order $J+b_d+1$ is left in
the energy bound. The word product convolution and the metric transport
cancellation consequently give a single radius-loss factor. In fact, in
a transport commutator with $1\le j\le n$, put $l=n-j+1$. The ratio of
Gevrey weights, including its Leibniz coefficient, is
\[
 \frac{r^n}{(n!)^2}\binom nj
   \frac{(j!)^2(l!)^2}{r^j r^l}
   =\frac{l^2}{r\binom nj}\le\frac lr,
\]
since $\binom nj\ge l$. The excluded term $j=0$ is precisely the
cancelled principal transport term. Thus the loss is the factor $l$
in $\mathcal L_J$, rather than an uncontrolled factor $l^2$ or an
extra derivative beyond the cutoff. We obtain
\begin{equation}\label{eq:packet-correction-energy}
 \mathcal A_J'
 \le C\{\mathcal A_J+\mathcal A_J^2+\varepsilon\}
  +\left\{\frac{r'}{r}
       +C(r^{-1}+R_c)(D_*/k+\mathcal A_J)\right\}\mathcal L_J.
\end{equation}
The same $C,R_c$ work for every external derivative cutoff $J$.
Here $\varepsilon=2\exp(-7X\log k/10)$ includes the metric conversion
of the residual after a fixed enlargement of constants.

Choose a positive $r_0$ from the known coefficient and packet radii so
that all pressure absorptions hold, in particular $r_0R_c\le1$, and set
\[
 \eta=\exp(-\sqrt X),\qquad
 r(t)=r_0-2C(D_*/k+\eta)t.
\]
If
\begin{equation}\label{eq:packet-correction-guards}
 2\varepsilon e^{3CT}\le\eta/2,
 \qquad 2C(D_*/k+\eta)T\le r_0/2,
\end{equation}
then $r_0/2\le r(t)\le r_0$. In the bootstrap range
$\mathcal A_J\le\eta\le1$ the loss coefficient in
\eqref{eq:packet-correction-energy} is nonpositive. The scalar integral
inequality gives the strict bound $\mathcal A_J\le\eta/2$ and closes
the bootstrap.

We give the existence step explicitly, following the pressure projection
and vanishing-viscosity argument of \cite[\S3.6]{OAI}. Let
$\mathscr G_\kappa$ be the closed gradient space above and let
$P_{\mathscr G}$ be its orthogonal projection in cylinder $L^2$.
On $\mathscr G_\kappa$ the map
$g\mapsto P_{\mathscr G}(G(t)g)$ is coercive, since
$\langle g,G(t)g\rangle\ge c\|g\|_2^2$ with $c>0$ supplied by $F$.
Set
\[
 \mathcal T_t=(P_{\mathscr G}G(t)|_{\mathscr G_\kappa})^{-1}
                         P_{\mathscr G},\qquad
 \mathcal P_t=I-G(t)\mathcal T_t.
\]
Then $P_{\mathscr G}\mathcal P_t=0$,
$\mathcal P_tu=u$ for $u\in\mathscr G_\kappa^\perp$, and
$\mathcal P_tG(t)g=0$ for $g\in\mathscr G_\kappa$.
Differentiating the coercive pressure equation gives the Sobolev bounds
for $\mathcal T_t$ and $\mathcal P_t$ at every fixed order. Their constants
may depend on that fixed order; the previously derived weighted energy
constants do not depend on the external derivative cutoff.

Write the unprojected correction source as
\[
 \mathcal N(E)=b_E\cdot\nabla_{y,\theta}E+\mathcal Z(E)+\mathcal E,
 \qquad b_E=\mathcal V_\kappa(Z+E),
\]
where $\mathcal Z(E)$ contains the terms $E\cdot\nabla_\kappa Z$,
$2F^{-1}F_tE$ and the differences of the two quadratic coefficient
terms. It has no derivative of $E$. For $0<\nu\le1$, solve
\begin{equation}\label{eq:packet-viscous-correction}
 E_t^\nu-\nu\Delta_{y,\theta}E^\nu
       =-\mathcal P_t\mathcal N(E^\nu),\qquad E^\nu(0)=0.
\end{equation}
The right side is locally Lipschitz $H^q\to H^{q-1}$ for
$q\ge b_d$, and the Fourier heat estimate
\[
 \|e^{\nu t\Delta}f\|_{H^q}
       \le C_q(1+(\nu t)^{-1/2})\|f\|_{H^{q-1}}
\]
gives a local mild solution by contraction. Every component of
$\nabla_\kappa$ has constant coefficients and commutes with the full
cylinder Laplacian. The heat semigroup therefore preserves the fixed
space $\mathscr G_\kappa^\perp$. Since the non-heat term in
\eqref{eq:packet-viscous-correction} takes values in this space, the mild
solution is exactly divergence free. No time derivative of
$\mathcal P_t$ and no commutation of $\mathcal P_t$ with the heat
semigroup is required.

The actual pressure gradient is
$-\mathcal T_t\mathcal N(E^\nu)$. The viscosity contributes no pressure,
because $\Delta E^\nu$ is still divergence free. Heat maximal regularity
justifies testing differentiated equations by
$F^{\top}F\partial^wE^\nu$: on each compact existence interval,
$E^\nu\in L^2_tH^{q+1}$ and $E_t^\nu\in L^2_tH^{q-1}$.
The differentiated metric-viscosity term satisfies
\[
 \nu\int\langle F^{\top}F u,\Delta u\rangle
       \le-\tfrac{c\nu}{2}\|\nabla u\|_2^2+C\nu\|u\|_2^2.
\]
This follows by integration by parts and Young's inequality, using the
known first derivatives of $F^{\top}F$. Thus it adds only a bounded
multiple of the lower energy to \eqref{eq:packet-correction-energy},
uniformly for $\nu\le1$. In the choice of $C$ above we include this
known metric-viscosity constant before selecting the radius or frequency;
it depends only on the fixed source and not on $J$ or $\nu$.
The same norm regularization justifies summing this estimate for
$\mathcal A_J$. The bootstrap extends every viscous solution
to $T$. Uniqueness in the low Sobolev class identifies the solutions
constructed at different higher orders. In particular they have uniform
bounds in $C_tH^q$ for each finite $q$.

There is no compactness assumption at spatial infinity in removing
viscosity. For two viscosities, subtract their equations and use the
same metric $L^2$ difference energy. The transport and pressure leading
terms cancel as before. The remaining viscous defect is
$(\nu-\mu)\Delta E^\mu$, uniformly bounded in $L^2$ by a constant
of the fixed source times $|\nu-\mu|$. With zero initial difference,
Gronwall gives
\[
 \|E^\nu-E^\mu\|_{C_tL^2}\le C_*|\nu-\mu|.
\]
The dissipative term $\nu\Delta(E^\nu-E^\mu)$ is handled by the same
integration by parts. Interpolation with the uniform higher-order
bounds now gives a strong limit in $C_tH^q$ for every finite $q$.
The nonlinear source and the actual pressure inverse pass to this limit
one order lower; $\nu\Delta E^\nu\to0$ there. The gradient and
solenoidal constraints are closed, the finite weighted energies pass to
the limit, and the initial trace remains zero. The limiting equation
supplies the continuous strong time derivative at every Sobolev order,
including the within-interval endpoint derivatives.

These steps use the heat Fourier multiplier on
$\R^d\times\mathbb T_P$, a closed Hilbert gradient space, and the
Sobolev product threshold \eqref{eq:packet-base-index}. In the
three-dimensional proof the corresponding cylinder has dimension four;
here the only changes are the $d+1$ Fourier variables, all $d$ quadratic
coefficient matrices, and the dimension-dependent Sobolev constants.
Neither a curl identity nor a two-dimensional transverse reduction
enters this existence argument.

There is an explicit finite frequency choice. With
\[
 X=k^{1/301000},\qquad N=\lfloor X\rfloor,
\]
\eqref{eq:packet-tail-base} meets $B_N\le k^{1/100}$ at sufficiently
large $k$. More concretely, after denoting by $I_*$ the inverse-multiplier
cost, it is enough for the packet and the correction guards that
\begin{equation}\label{eq:packet-explicit-threshold}
 k\ge\left(64+e+2^{300}C_*+I_*+12CT
       +\frac{8CTD_*}{r_0}
       +\left(\frac{8CT}{r_0}\right)^2\right)^{301000}.
\end{equation}
Indeed the $301000$th root controls $X$, the polynomial tail, $\log k$,
and the two terms in the radius loss; the negative exponential in
\eqref{eq:packet-small-residual} dominates $e^{3CT}$ and
$e^{-\sqrt X}$. The constants in \eqref{eq:packet-explicit-threshold}
are quantities of the current known source. This statement gives no
bound for them in terms of the parameters of a future parent.

\subsection{The physical child and the signed pressure term}
Define the child by
\begin{equation}\label{eq:packet-exact-child}
 U^+(t,x)=U(t,x)+\kappa F(t,y)
          W(t,y,km_0\cdot y),\qquad y=\Phi_t^{-1}(x).
\end{equation}
The chain rule, $\Phi_t'=U\circ\Phi_t$, and
\eqref{eq:packet-frame} transform
\eqref{eq:packet-coordinate-equation} into the physical momentum equation
with pressure force
\[
 \nabla P_U(t,x)+\kappa F(t,y)^{-\top}
             \mathbf p(t,y,km_0\cdot y).
\]
This force is an actual gradient on physical space. To see this without
assuming a cylinder scalar potential, observe that closed-gradient
membership implies
\[
 D_{\kappa,i}\mathbf p_j-D_{\kappa,j}\mathbf p_i=0
\]
in distributions: it holds for each compact smooth gradient generator
and is preserved by the $L^2$ limit. All spatial derivatives of the
constructed pressure vector are continuous, so the same identity holds
pointwise. The pullback chain rule for a one-form gives
\[
 \partial_{x_i}(\kappa F^{-\top}\mathbf p)_j
 -\partial_{x_j}(\kappa F^{-\top}\mathbf p)_i
 =\sum_{a,b}(F^{-1})_{bi}(F^{-1})_{aj}
       (D_{\kappa,b}\mathbf p_a-D_{\kappa,a}\mathbf p_b)=0
\]
after graph evaluation. The terms differentiating $F^{-1}=D\Phi_t^{-1}$
cancel by the symmetry of its ordinary second derivatives.

Graph evaluation uses the smooth periodic representatives and one extra
angular Sobolev trace estimate; the parent change of variables is volume
preserving. It gives all physical $L^2$ spatial derivatives and their
continuous time paths, both for the velocity and for this pressure force.
A curl-free $L^2$ vector field on $\R^d$ lies in the ordinary closed
gradient space, by its Fourier characterization. The actual Piola
identity also transforms the closed cylinder solenoidal constraint into
$\operatorname{div}U^+=0$. Taking the ordinary divergence of the physical
momentum equation and applying the Helmholtz decomposition identifies
its pressure force with the gradient of the Newton pressure in
Lemma~\ref{lem:newton}. We choose exactly that normalized scalar pressure.
The smooth Sobolev local theory then identifies the constructed field
with the actual Euler evolution on the whole interval. Auxiliary mean
and finite angular scalars fixed convenient representatives, but no
global periodic scalar for the correction was required.

The same field has the decomposition
\[
 U^+=U+V^{[N]}(t,y,km_0\cdot y)
        +\kappa F(t,y)E(t,y,km_0\cdot y).
\]
Thus its initial increment is exactly the finite packet, and the
correction initially vanishes. The all-order correction estimate gives
physical velocity-gradient and pressure-Hessian errors bounded by fixed
current-source constants times $\eta$. One fast derivative of the velocity correction is paired with its
prefactor $\kappa$. The physical pressure force of the correction is
$\kappa F^{-\top}\mathbf p^e$ evaluated on the graph. Its one
physical derivative, which is the pressure Hessian, has the same
cancellation of $\kappa$ against one fast derivative. For the finite
scalar primary pressure the equivalent accounting is its prefactor
$\kappa^2$ against two fast derivatives. All remaining powers of
$\kappa$ are nonnegative. A fixed number of the uniform weighted
estimates controls the required coefficient and angular derivatives;
the stated constants are independent of $k$.

The periodic profile in \eqref{eq:packet-joined-primary} can be chosen
with explicit polynomial dependence on its width. Fix an even Gevrey--2
function $\eta$ supported in $[-1,1]$, with $0\le\eta\le1$ and
$\eta(0)=1$, and write $m_\eta=\int_\R\eta>0$.
For example, take $\eta(s)=\exp(1-(1-s^2)^{-1})$ on $|s|<1$
and zero elsewhere. On a complex disc of radius $c(1-s^2)$ about
a real $s\in(-1,1)$, its modulus is at most
$C\exp(-c'/(1-s^2))$. Cauchy's estimate, followed by maximizing
$t^{-n}e^{-c'/t}$, yields $\|\eta^{(n)}\|_\infty\le C R^n(n!)^2$;
the same estimate gives the zero derivative traces at $\pm1$.
Choose a fixed $a>0$ with $a\le P/4$ and $am_\eta\le P/4$.
For $0<\delta\le1$, let $w=a\delta$, let $\eta_w$ denote the
$P$-periodization of $\eta(\theta/w)$, and put $q_w=wm_\eta/P$.
Define $f_\delta$ by
\[
 f_\delta(0)=0,\qquad
 f_\delta'(\theta)=\frac{\eta_w(\theta)-q_w}
                            {\delta(1-q_w)}.
\]
The derivative is even and has zero angular mean, so this defines an
odd $P$-periodic function with zero mean. Since $0<q_w\le1/4$,
\[
 f_\delta'(0)=\delta^{-1},\qquad
 -\frac13\le f_\delta'\le\delta^{-1}.
\]
Indeed $q_w/\delta=am_\eta/P\le1/4$. If
$\|\eta^{(n)}\|_\infty\le C_0R_0^n(n!)^2$, differentiating this
formula gives, for constants $C_*,R_*$ depending only on the fixed bump
and period,
\[
 \|f_\delta^{(n)}\|_\infty
       \le C_*(R_*/\delta)^n(n!)^2\quad(n\ge0).
\]
At order zero, integrate $|f_\delta'|$ over a period: its positive and
negative parts have equal integrals and each is bounded by a constant
independent of $\delta$. For $n\ge1$, the factors are
$\delta^{-1}w^{-(n-1)}$ and the derivatives of the fixed bump.
Thus the profile's Gevrey radius is a polynomial source cost.

Use this profile for the joined primary \eqref{eq:packet-joined-primary}.
The uncut physical high velocity is $b_Y$, and the amplitude is $\alpha$.
At
$y=\Phi_t^{-1}(x)$ and $\vartheta=km_0\cdot y$, the leading velocity
and pressure Hessian tensors are respectively
\begin{align}
 \mathcal S&=\alpha\chi_{\rm in}(y)f_\delta'(\vartheta)\,
                   b_Y\otimes m,\notag\\
 \mathcal Q&=-2\alpha\chi_{\rm in}(y)
       \frac{\langle m,Mb_Y\rangle}{|m|^2}
       f_\delta'(\vartheta)\,m\otimes m.
                                                        \label{eq:packet-leading-tensors}
\end{align}
The actual finite higher grades, derivatives of slow coefficients, and
exact correction give
\begin{equation}\label{eq:packet-low-errors}
 \|DU^+-DU-\mathcal S\|\le C_v/k+C'_v\eta,
 \quad
 \|D^2P_{U^+}-D^2P_U-\mathcal Q\|\le C_p/k+C'_p\eta,
\end{equation}
uniformly on the closed time interval and physical space. The pressure
identity follows first for the true finite scalar pressure and then for
the true correction pressure; equality of their gradient with the
canonical Euler pressure gradient transfers the identity to
$P_{U^+}$.

If $\langle m,Mb_Y\rangle\ge0$, the lower slope bound gives the one-sided
estimate
\begin{equation}\label{eq:packet-signed-curvature}
 \langle\mathcal Qz,z\rangle
 \le2\alpha\chi_{\rm in}\langle m,Mb_Y\rangle |z|^2.
\end{equation}
Without that sign the valid absolute estimate is
\[
 \|\mathcal Q\|\le
   2\alpha\|M\|\chi_{\rm in}|m|\,|b_Y|/\delta.
\]
The distinction is necessary: using the absolute bound at late times
would lose the narrow-profile gain in the next mean curvature budget.
All tensors in \eqref{eq:packet-leading-tensors} act on the entire
$\R^d$, and the norm identities used here hold in every dimension.

For spatial localization the preceding construction is applied to the
actual rescaled parent
\[
 U^\rho(t,x)=\rho^{-1}U(t,\rho x),\qquad
 P_{U^\rho}(t,x)=\rho^{-2}P_U(t,\rho x),\qquad 0<\rho\le1.
\]
Direct substitution proves the Euler equations on the same time interval.
Its actual flow is $\rho^{-1}\Phi_t(\rho y)$, so its frame, strain and
pressure Hessian in labels are exactly $F(t,\rho y)$,
$M(t,\rho y)$ and $H(t,\rho y)$. Undoing the scaling after constructing
the child gives the literal physical field
$\rho(U^\rho)^+(t,x/\rho)$. Velocity gradients and pressure Hessians
are unchanged in size by these two inverse scalings, and the support
radius of a compact oscillatory coefficient is multiplied by $\rho$.
Thus the low error bounds and leading tensors above transfer to the
original parent with no factor $\rho^{-1}$ in those two derivative
orders. Higher-order initial derivatives retain their true scaling
factors as recorded in Section~\ref{sec:initial-summability}.

\subsection{Physical localization of the mean boundary condition}
The oscillatory localization scale and the mean cutoff scale are chosen
separately. With a fixed physical cutoff parameter $\ell_*$ and physical
unscaling $x=\rho y$, use the normalized cutoff parameter
$\ell=\rho\ell_*$. Then
$\chi(\ell y)=\chi(\ell_*x)$, and
\eqref{eq:packet-mean-initial} places each physical mean initial
increment in $\overline B(0,\ell_*^{-1})$. If
$\operatorname{supp}\chi_{\rm in}\subset\overline B(0,1/4)$, the high
and corrector initial increments are supported in
$\overline B(0,\rho/4)$. The exact correction has initial value zero.
Lemma~\ref{lem:fixed-cutoff} proves propagation of the mean reserve
with this fixed physical cutoff, and Lemma~\ref{lem:common-support}
gives the resulting common compact support through the iteration.


\subsection{The one-parent packet conclusion}
We collect the output needed by the iteration. The assumptions in the
following lemma concern only the current, already known source.

\begin{lemma}[Exact joined packet child]\label{lem:exact-packet-child}
Fix $d\ge3$ and a smooth Sobolev parent on $[0,T]$ with actual flow
\eqref{eq:packet-frame}, the known coefficient bounds
\eqref{eq:packet-known-jets}, and the mean reserve
\eqref{eq:packet-mean-guard}. Fix its transverse terminal datum, compact
Gevrey cutoffs, and the positive forward profile with
\eqref{eq:packet-tail-propagator}. There is a finite constant determined
by these data such that every sufficiently large $k$ gives the actual
child \eqref{eq:packet-exact-child} from the joined strict-prefix family.
It solves the full $d$-dimensional Euler equation with its Newton pressure
on the whole closed interval and belongs to
$C_tH^m\cap C_t^1H^{m-1}$ for every integer $m\ge1$.

Its leading velocity-gradient and pressure-Hessian increments are
exactly the two tensors in \eqref{eq:packet-leading-tensors}. After an
explicit enlargement of the known-source frequency threshold, each
error in \eqref{eq:packet-low-errors} is at most $k^{-1/4}$, uniformly in
time and space. Its initial increment is the entire finite sum
\eqref{eq:packet-finite-pairs} at time zero, including every higher high
initial velocity and the terminal corrector $C_N$; the exact correction
has initial value zero.

If the parent is odd and the cutoffs are even, then the child is odd.
If in addition a fixed orthogonal coordinate reflection $\mathscr R$
obeys
\[
 U(t,\mathscr Rx)=\mathscr R U(t,x),\quad
 \mathscr R m_0=m_0,\quad \mathscr R R_\perp Y=R_\perp Y,
\]
and preserves the cutoffs, then
\[
 U^+(t,\mathscr Rx)=\mathscr R U^+(t,x),\qquad
 DU^+(t,\mathscr Rx)=\mathscr R DU^+(t,x)\mathscr R^{-1}.
\]
For the fixed physical mean cutoff described above, every
initial increment is supported in
$\overline B(0,\max(1/4,\ell_*^{-1}))$. Thus the child's initial support
is contained in the union of this ball and the parent's initial support.
\end{lemma}
\begin{proof}
The constructions and identities above give existence, the actual PDE,
the normalized pressure, and all Sobolev orders. The terminal corrector
and pressure were included before computing the literal residual, so
no additional endpoint grade is omitted.

Here is a concrete enlargement giving the stated error. Let
$\beta=301000$ and let $C_{\rm err}\ge1$ dominate the four constants in
\eqref{eq:packet-low-errors}. Increase the lower bound for $X=k^{1/\beta}$
to include $(2\beta)^4$ and $(4C_{\rm err})^2$, as well as the quantity
inside parentheses in \eqref{eq:packet-explicit-threshold}.
For $X\ge1$, $\log X\le4X^{1/4}$, and hence
\[
 \tfrac14\log k\le\tfrac12\sqrt X,
 \qquad \log(4C_{\rm err})\le\tfrac12\sqrt X.
\]
It follows that
$C_{\rm err}e^{-\sqrt X}\le k^{-1/4}/4$ and
$C_{\rm err}/k\le k^{-1/4}/4$. The sum satisfies the asserted error
bound. The enlarged threshold remains a fixed power of a sum of known
source constants.

For symmetry, uniqueness of the parent's flow first gives
$\Phi_t(-y)=-\Phi_t(y)$ and the corresponding reflection covariance.
The mean form commutes with the induced orthogonal action, since its
cutoff tensor operator is constructed using all tensor indices and the
cutoff is invariant. Uniqueness of its coercive inverse gives the mean
parities. For the joined high operator the action preserves the fixed
zero-trace displacement space; for the primary it also preserves the
terminal datum. Uniqueness of the history inverse and then of the
forward initial-value equation proves the high parities. The normalized
angular primitive, the full tensor corrector, and each literal product
in the strict-prefix force commute with the same actions. Induction
therefore gives odd vector fields and even scalar pressures for the
whole finite family, together with the stated reflection covariance.

The projected coordinate equation and its closed solenoidal and gradient
spaces are also invariant under these actions. Its unique correction
with zero initial value has the same parities. Transforming back gives
the actual child symmetry; differentiating this identity gives the
claimed full derivative covariance. Finally, the initial support
follows from \eqref{eq:packet-mean-initial}, support preservation of the
joined high inverse and the tensor potential, and $E(0)=0$.
\end{proof}
